\section{Conclusion and Future Horizons}
\label{sec:conclusion}

In this paper, we challenged the foundational assumption of model merging---that combined task experts must reside as a static, input-independent point in parameter space. Inspired by quantum mechanics, we proposed \textbf{Quantum Wavefunction Superposition Merging (QWS-Merge)}, a radical new paradigm that models individual model weights as eigenstates in a parameter Hilbert space. Under this formulation, merging is represented as a coherent, quantum-like superposition of parameter states that dynamically collapses to a localized weight configuration based on the wave-like phase interference of incoming input features.

Our experimental evaluation on a challenging heterogeneous visual benchmark (MNIST, FashionMNIST, CIFAR-10, SVHN) using a compact Vision Transformer backbone demonstrated that standard static and unsupervised test-time merging techniques suffer from catastrophic representational collapse. In stark contrast, QWS-Merge completely resolves parameter sign and gradient conflicts by dynamically routing activations through optimized task-specific parameter manifolds, elevating the multi-task average accuracy from $49.35\%$ (Uniform Merging) to \textbf{59.32\%}. Crucially, when compared against an unconstrained classical Linear Router, QWS-Merge exhibits strong wave-inspired regularization: on the highly conflicting SVHN dataset, the Linear Router collapses to $15.30\%$, while QWS-Merge preserves $91.5\%$ of specialized expert capacity at \textbf{31.60\%}. Additionally, we provide a transparent, rigorous study of batch size and task heterogeneity on dynamic routers, shedding light on the "heterogeneity collapse" at larger batch sizes.

We believe QWS-Merge represents a significant, non-incremental milestone that shifts the conversation from static parameter blending to dynamic parameter-wave routing. Future horizons include scaling this quantum-inspired paradigm to large-scale language models and multimodal vision-language architectures, where task conflicts are diverse and complex. Additionally, exploring adaptive wave frequencies and complex-valued parameters could unlock richer wave interference dynamics, further expanding the capabilities of parameter-space multi-task learning.
